\documentclass[10pt,aspectratio=169]{beamer}
\input{../../_en_preambulo_beamer}
% Comandos y macros estadísticas compartidas para las presentaciones Beamer en inglés (Metropolis)

% Conjuntos numéricos básicos
\newcommand{\R}{\mathbb{R}}
\newcommand{\N}{\mathbb{N}}
\newcommand{\Q}{\mathbb{Q}}
\newcommand{\Z}{\mathbb{Z}}
\newcommand{\set}[1]{\left\{ #1 \right\}}
\newcommand{\abs}[1]{\left|#1\right|}
\newcommand{\norm}[1]{\left\| #1 \right\|}

% Comandos estadísticos y probabilísticos del curso
\newcommand{\Var}{\operatorname{Var}}
\newcommand{\cov}{\operatorname{Cov}}
\newcommand{\corr}{\rho}
\newcommand{\comb}[2]{\begin{pmatrix} #1 \\ #2 \end{pmatrix}}
\newcommand{\s}{\sigma}
\newcommand{\particion}{\mathfrak{P}}
\newcommand{\card}[1]{\#\left( #1 \right)}
\newcommand{\seq}[3]{\set{#1}_{#2}^{#3}}
\newcommand{\E}{\mathbb{E}}
\newcommand{\Prob}{\mathbb{P}}

% Entornos pedagógicos y cajas en inglés académico natural (soportando tanto nombres en inglés como en español)
\newenvironment{discussion}{%
  \begin{alertblock}{Discussion Question}%
}{%
  \end{alertblock}%
}
\newenvironment{discusion}{%
  \begin{alertblock}{Discussion Question}%
}{%
  \end{alertblock}%
}

\newenvironment{reflection}{%
  \begin{alertblock}{Classroom Reflection}%
}{%
  \end{alertblock}%
}
\newenvironment{reflexion}{%
  \begin{alertblock}{Classroom Reflection}%
}{%
  \end{alertblock}%
}

\newenvironment{paradox}{%
  \begin{alertblock}{Exploratory Paradox}%
}{%
  \end{alertblock}%
}
\newenvironment{paradoja}{%
  \begin{alertblock}{Exploratory Paradox}%
}{%
  \end{alertblock}%
}

\newenvironment{motivation}{%
  \begin{exampleblock}{Motivation: Data Science in Practice}%
}{%
  \end{exampleblock}%
}
\newenvironment{motivacion}{%
  \begin{exampleblock}{Motivation: Data Science in Practice}%
}{%
  \end{exampleblock}%
}

\newenvironment{quickchallenge}{%
  \begin{block}{Quick Team Challenge}%
}{%
  \end{block}%
}
\newenvironment{retoclase}{%
  \begin{block}{Quick Team Challenge}%
}{%
  \end{block}%
}


\title[Introduction to Inferential Statistics]{Section 5.0: Introduction to Inferential Statistics}
\subtitle{From a sample to population conclusions}
\author[J. Castillo Colmenares]{Juliho Castillo Colmenares}
\institute{Tecnológico de Monterrey}
\date{}

\begin{document}
\begin{frame}[plain]\vspace{-0.8cm}\titlepage\end{frame}

\begin{frame}{Roadmap --- Chapter 5}
  \footnotesize
  \begin{itemize}\itemsep=0.08cm
    \item \textbf{05.00} Introduction to inferential statistics \emph{(today)}
    \item \textbf{05.01--05.02} Transformations and sampling distributions
    \item \textbf{05.03--05.05} Chi-square, $t$, and $F$
    \item \textbf{05.08--05.09} Fundamental concepts and Z/t statistics
  \end{itemize}
  \begin{block}{Learning objective}
    Distinguish population, sample, parameter, and statistic; describe the
    inference process and the role of sampling error and the CLT.
  \end{block}
\end{frame}

\begin{frame}{Why infer?}
  \footnotesize
  The full population is rarely observable. The reading identifies four
  practical constraints:
  \begin{itemize}\itemsep=0.1cm
    \item cost of measuring everyone;
    \item collection time;
    \item limited accessibility;
    \item destructive measurements, such as testing a lightbulb's lifetime.
  \end{itemize}
  \begin{alertblock}{Central idea}
    A representative sample supports generalization, but uncertainty must be
    quantified.
  \end{alertblock}
\end{frame}

\begin{frame}{Population and sample}
  \footnotesize
  \begin{block}{Definitions from the section}
    \textbf{Population:} complete set of individuals, objects, or events with a
    common characteristic of interest.

    \medskip
    \textbf{Sample:} subset selected for the study.
  \end{block}
  \begin{exampleblock}{Operational reading}
    The population sets the scope of the conclusion; the sample supplies the
    observed data.
  \end{exampleblock}
\end{frame}

\begin{frame}{Parameter and statistic}
  \footnotesize
  \begin{columns}[T]
    \column{0.48\textwidth}
      \begin{block}{Population}
        Fixed, usually unknown parameter: $\mu$, $\sigma$.
      \end{block}
    \column{0.48\textwidth}
      \begin{block}{Sample}
        Computable statistic: $\bar{x}$, $s$.
      \end{block}
  \end{columns}
  \begin{alertblock}{Distinction}
    A parameter describes the population; a statistic changes from sample to
    sample and is used to estimate it.
  \end{alertblock}
\end{frame}

\begin{frame}{Two forms of inference}
  \footnotesize
  \begin{enumerate}\itemsep=0.18cm
    \item \textbf{Estimation:} use statistics to approximate a parameter,
      either pointwise or through a confidence interval.
    \item \textbf{Hypothesis testing:} state a claim about a parameter and
      decide whether sample evidence supports or contradicts it.
  \end{enumerate}
  \begin{exampleblock}{Guiding question}
    Do we need a plausible value/range, or a decision about a specific claim?
  \end{exampleblock}
\end{frame}

\begin{frame}{General inference process}
  \footnotesize
  \begin{enumerate}\itemsep=0.12cm
    \item Define the population and characteristic of interest.
    \item Select a representative sample.
    \item Calculate sample statistics.
    \item Apply estimation or testing with suitable distributions.
    \item Interpret the result in context.
  \end{enumerate}
  \begin{block}{Logical sequence}
    \centering population $\longrightarrow$ sample $\longrightarrow$ statistic
    $\longrightarrow$ inference $\longrightarrow$ decision.
  \end{block}
\end{frame}

\begin{frame}{Sources of error}
  \footnotesize
  \begin{columns}[T]
    \column{0.48\textwidth}
      \begin{block}{Sampling error}
        Random difference between statistic and parameter; it decreases as $n$
        grows.
      \end{block}
    \column{0.48\textwidth}
      \begin{block}{Non-sampling error}
        Selection, measurement, or non-response bias; it does not vanish just
        by increasing $n$.
      \end{block}
  \end{columns}
  \begin{alertblock}{Warning}
    Inference quantifies mainly sampling error; study design must control
    systematic errors.
  \end{alertblock}
\end{frame}

\begin{frame}{Sampling distribution}
  \footnotesize
  If many samples of size $n$ are taken and a statistic is calculated for each,
  the results form its \emph{sampling distribution}.
  \begin{block}{Conceptual example from the reading}
    Each sample produces a mean $\bar X$. The distribution of all means measures
    estimation variability, not individual-observation variability.
  \end{block}
  \begin{alertblock}{Bridge}
    This distribution underlies confidence intervals and hypothesis tests.
  \end{alertblock}
\end{frame}

\begin{frame}{Computational bridge 1/4: quantities}
  \footnotesize
  \begin{tabular}{lll}
    \hline
    Object & Symbol & Role in inference \\
    \hline
    Population & $\mu,\sigma$ & target values \\
    Sample & $X_1,\ldots,X_n$ & observed data \\
    Statistic & $\bar X,s$ & computable quantities \\
    Inference & CI or test & uncertain conclusion \\
    \hline
  \end{tabular}
\end{frame}

\begin{frame}{Computational bridge 2/4: sample mean}
  \footnotesize
  For a random sample with population mean $\mu$ and variance $\sigma^2$,
  \[
    \bar X=\frac{1}{n}\sum_{i=1}^{n}X_i.
  \]
  The sample mean connects observed data with the target parameter $\mu$.
\end{frame}

\begin{frame}{Computational bridge 3/4: variability}
  \footnotesize
  \[
    E(\bar X)=\mu,\qquad \operatorname{Var}(\bar X)=\frac{\sigma^2}{n}.
  \]
  \begin{itemize}\itemsep=0.13cm
    \item The mean is unbiased.
    \item Its standard error is $\sigma/\sqrt n$.
    \item Quadrupling $n$ halves the standard deviation of the mean.
  \end{itemize}
\end{frame}

\begin{frame}{Computational bridge 4/4: CLT}
  \footnotesize
  \[
    \frac{\bar X-\mu}{\sigma/\sqrt n}\xrightarrow{d}N(0,1)
    \qquad(n\text{ large}).
  \]
  The original population shape becomes less important for a sufficiently large
  sample mean.
\end{frame}

\begin{frame}{Reading case 1/4 --- population and sample}
  \footnotesize
  A destructive measurement prevents studying every object in a population.
  The reading proposes a representative sample.
  \begin{alertblock}{Conceptual decision}
    The sample does not replace the population definition; it provides the data
    used to learn about it.
  \end{alertblock}
\end{frame}

\begin{frame}{Reading case 1/4 --- interpretation}
  \footnotesize
  \begin{block}{Conceptual solution}
    Set the target population, select the sample, and quantify the uncertainty
    of generalizing from one to the other.
  \end{block}
  A larger sample can reduce sampling error, but not selection or measurement
  bias by itself.
\end{frame}

\begin{frame}{Reading case 2/4 --- parameter and statistic}
  \footnotesize
  The population mean $\mu$ is fixed and unknown. The sample mean $\bar X$
  changes when sampling is repeated.
  \begin{block}{Question}
    Which quantity can be calculated directly from the data?
  \end{block}
  The section's answer is the statistic $\bar X$, used to estimate $\mu$.
\end{frame}

\begin{frame}{Reading case 2/4 --- estimation}
  \footnotesize
  \[
    \widehat{\mu}=\bar X.
  \]
  \begin{itemize}\itemsep=0.12cm
    \item Point estimation gives one value.
    \item Interval estimation adds a range and confidence level.
    \item Both depend on the sampling distribution.
  \end{itemize}
\end{frame}

\begin{frame}{Reading case 3/4 --- inference process}
  \footnotesize
  \begin{enumerate}\itemsep=0.13cm
    \item Define the characteristic.
    \item Select the sample.
    \item Calculate the statistic.
    \item Choose the inferential procedure.
    \item Interpret in context.
  \end{enumerate}
\end{frame}

\begin{frame}{Reading case 3/4 --- sources of error}
  \footnotesize
  \begin{block}{Conceptual solution}
    Sampling error is modeled through the sampling distribution; non-sampling
    errors require design, measurement, and follow-up.
  \end{block}
  \begin{alertblock}{Conclusion}
    A valid inference addresses both types of error.
  \end{alertblock}
\end{frame}

\begin{frame}{Reading case 4/4 --- role of the CLT}
  \footnotesize
  Under its conditions and for sufficiently large $n$, the CLT approximates the
  sampling distribution of a mean by a normal distribution.
  \begin{block}{Later use}
    This approximation supports Z/t tests and confidence intervals developed in
    the chapter.
  \end{block}
\end{frame}

\begin{frame}{Synthesis of the section}
  \footnotesize
  \begin{itemize}\itemsep=0.12cm
    \item Population and sample set the scope of the problem.
    \item Parameters and statistics distinguish target from evidence.
    \item Estimation and testing are the two inferential forms.
    \item Sampling error and the CLT explain uncertainty.
  \end{itemize}
\end{frame}

\begin{frame}{Curricular transition}
  \footnotesize
  \begin{block}{Established}
    Inference moves from sample statistics to population claims while recording
    error and assumptions.
  \end{block}
  \begin{alertblock}{Next section}
    \textbf{05.01 Transformation of Variables}: obtain the distribution of
    $Y=g(X)$ and prepare the study of sampling distributions.
  \end{alertblock}
\end{frame}
\end{document}
